\chapter[1967 --- Granit et al.]{1967: Ragnar Granit, Haldan Keffer Hartline, George Wald}
\label{ch:1967_Granit}

\section*{Main Contribution}

\textit{Discovered the primary physiological and chemical processes of vision, explaining how light is converted into nerve signals in the retina.}

\section{Official Citation}

for their discoveries concerning the primary physiological and chemical visual processes in the eye

\section{Background and Historical Context}

The 1967 Nobel Prize in Physiology or Medicine was awarded to Ragnar Granit, Haldan Keffer Hartline, and George Wald for their discoveries concerning the primary physiological and chemical visual processes in the eye. This prize recognized a series of breakthroughs that collectively explained how the eye converts light into neural signals---a process known as phototransduction.

\subsection{The Mystery of Vision}

The eye has fascinated scientists for centuries. By the mid-20th century, it was well established that the eye contains photoreceptor cells (rods and cones) that respond to light and generate electrical signals that are transmitted to the brain. However, the mechanisms by which light is converted into neural signals were still poorly understood.

The work of Granit, Hartline, and Wald would provide the first detailed understanding of the visual process, from the absorption of light by visual pigments to the generation of electrical signals in the retina.

\subsection{Ragnar Granit}

Ragnar Granit was born on October 30, 1900, in Riihim\`{a}ki, Finland. He studied at the University of Helsinki, where he earned his doctorate in 1927. After completing his doctorate, Granit worked at the University of Oxford and later at the Karolinska Institute in Stockholm, where he would spend most of his career.

Granit was a skilled physiologist who was interested in the electrical activity of the retina. His work on the electroretinogram (ERG)---the electrical response of the retina to light---would earn him the Nobel Prize.

\subsection{Haldan Keffer Hartline}

Haldan Keffer Hartline was born on December 22, 1903, in Bloomsburg, Pennsylvania. He studied at Lafayette College, where he earned his bachelor's degree in 1923, and at the Johns Hopkins University, where he earned his medical degree in 1927. After completing his medical training, Hartline worked at the Rockefeller Institute for Medical Research and later at the University of Pennsylvania, where he would spend most of his career.

Hartline was a skilled neurophysiologist who was interested in the electrical activity of individual nerve cells in the retina. His work on the electrical response of single nerve fibers to light would earn him the Nobel Prize.

\subsection{George Wald}

George Wald was born on November 18, 1906, in New York City. He studied at New York University, where he earned his bachelor's degree in 1927, and at Columbia University, where he earned his Ph.D. in 1932. After completing his doctorate, Wald worked at the University of Chicago and later at Harvard University, where he would spend most of his career.

Wald was a skilled biochemist who was interested in the chemistry of the visual pigments. His discovery of the chemical structure of visual pigments would earn him the Nobel Prize.

\subsection{The Three Components of Vision}

The visual process involves three main components:

\begin{itemize}
\item The absorption of light by visual pigments in the photoreceptor cells
\item The conversion of light energy into electrical signals in the photoreceptor cells
\item The transmission of these electrical signals to the brain through the optic nerve
\end{itemize}

The work of Granit, Hartline, and Wald addressed each of these components, providing a comprehensive understanding of the visual process.

\section{The Discovery/Contribution}

The work of Granit, Hartline, and Wald on visual processes involved the discovery of the mechanisms by which light is absorbed, converted into electrical signals, and transmitted to the brain.

\subsection{Granit's Work on the Electroretinogram}

Ragnar Granit's primary contribution was his work on the electroretinogram (ERG)---the electrical response of the retina to light.

\subsubsection{The Discovery}

Granit used a sensitive galvanometer to measure the electrical activity of the retina in response to light. He discovered that the retina produces a characteristic electrical response when stimulated by light, which he called the electroretinogram (ERG).

\subsubsection{Analysis of the ERG}

Granit analyzed the ERG and showed that it consists of several components:

\begin{itemize}
\item The a-wave: An initial negative deflection that reflects the activity of the photoreceptor cells (rods and cones)
\item The b-wave: A subsequent positive deflection that reflects the activity of the bipolar cells
\item The c-wave: A slow positive deflection that reflects the activity of the retinal pigment epithelium
\end{itemize}

\subsubsection{Significance of the Discovery}

Granit's work on the ERG was significant because it provided a non-invasive method for studying the electrical activity of the retina. The ERG is now widely used in clinical ophthalmology to diagnose and monitor retinal diseases, such as retinitis pigmentosa and diabetic retinopathy.

\subsection{Hartline's Work on Single Nerve Fibers}

Haldan Keffer Hartline's primary contribution was his work on the electrical response of single nerve fibers in the retina to light.

\subsubsection{The Discovery}

Hartline used microelectrodes to record the electrical activity of individual nerve fibers in the retina of the horseshoe crab (\textit{Limulus polyphemus}). He discovered that individual nerve fibers respond to light by generating electrical impulses (action potentials) whose frequency increases with the intensity of the light.

\subsubsection{Lateral Inhibition}

Hartline also discovered a phenomenon called lateral inhibition, in which the response of one photoreceptor cell inhibits the response of its neighbors. This process enhances the contrast of visual images by emphasizing the edges and boundaries of objects.

\subsubsection{Significance of the Discovery}

Hartline's work on single nerve fibers was significant because it provided the first detailed understanding of how individual nerve cells in the retina process visual information. His discovery of lateral inhibition explained how the retina enhances the contrast of visual images, a process that is essential for visual perception.

\subsection{Wald's Work on Visual Pigments}

George Wald's primary contribution was the discovery of the chemical structure of visual pigments and the elucidation of the visual cycle.

\subsubsection{The Discovery}

Wald discovered that the visual pigments in the retina are composed of a protein called opsin and a light-absorbing molecule called retinal. Retinal is derived from vitamin A, which is why vitamin A deficiency can cause night blindness.

\subsubsection{The Visual Cycle}

Wald elucidated the visual cycle, the process by which visual pigments are regenerated after being bleached by light. The visual cycle involves the following steps:

\begin{itemize}
\item Light is absorbed by rhodopsin, the visual pigment in rods, causing retinal to change shape from 11-cis to all-trans
\item This conformational change triggers a series of chemical reactions that ultimately generate an electrical signal in the rod cell
\item The all-trans retinal is released from opsin and is converted back to 11-cis retinal by a series of enzymatic reactions
\item The 11-cis retinal recombines with opsin to form rhodopsin, which is ready to absorb light again
\end{itemize}

\subsubsection{Significance of the Discovery}

Wald's work on visual pigments was significant because it provided the first detailed understanding of the chemical basis of vision. His discovery that vitamin A is a precursor to retinal explained why vitamin A deficiency causes night blindness and led to the development of treatments for this condition.

\section{Impact on Modern Medicine}

The work of Granit, Hartline, and Wald on visual processes has had a significant impact on modern medicine. Their discoveries have contributed to the understanding and treatment of eye diseases and have led to the development of numerous medical technologies and therapies.

\subsection{Diagnosis of Eye Diseases}

The ERG developed by Granit is now widely used in clinical ophthalmology to diagnose and monitor retinal diseases. The ERG can detect early signs of retinal damage before symptoms appear, allowing for earlier diagnosis and treatment of eye diseases.

\subsection{Understanding Night Blindness}

Wald's work on visual pigments has contributed to our understanding of night blindness and other visual disorders caused by vitamin A deficiency. The discovery that vitamin A is a precursor to retinal has led to the development of treatments for night blindness, including vitamin A supplementation.

\subsection{Age-Related Macular Degeneration}

The understanding of visual processes that Granit, Hartline, and Wald helped to establish has also contributed to our understanding of age-related macular degeneration (AMD), a common eye disease that causes vision loss in older adults. Research into the mechanisms of visual pigments has led to the development of new treatments for AMD, including anti-VEGF (vascular endothelial growth factor) drugs.

\subsection{Retinal Prostheses}

The understanding of the visual process has also contributed to the development of retinal prostheses---electronic devices that can restore some vision to patients with retinal diseases. These devices work by electrically stimulating the remaining retinal cells, mimicking the natural process of phototransduction.

\subsection{Gene Therapy for Eye Diseases}

The understanding of the genetic basis of eye diseases has also contributed to the development of gene therapy for these conditions. Gene therapy involves introducing functional copies of defective genes into the retina, restoring normal function. This approach has shown promise for treating a variety of genetic eye diseases, including Leber congenital amaurosis and retinitis pigmentosa.

\section{Legacy}

The legacy of Ragnar Granit, Haldan Keffer Hartline, and George Wald extends far beyond their Nobel Prize-winning work on visual processes. Their discoveries laid the foundation for modern visual science and have contributed to numerous advances in the diagnosis and treatment of eye diseases.

\subsection{Foundation for Visual Science}

The work of Granit, Hartline, and Wald established the foundations of modern visual science. Their discoveries about the electroretinogram, single nerve fiber responses, and visual pigments provided a comprehensive understanding of the visual process that continues to guide research in visual science today.

\subsection{Advances in Ophthalmology}

The understanding of visual processes has contributed to numerous advances in ophthalmology, including the development of new diagnostic tests and treatments for eye diseases. These advances have improved the lives of millions of patients with eye diseases and continue to drive progress in this field.

\subsection{Inspiration for Future Researchers}

The work of Granit, Hartline, and Wald continues to inspire researchers in visual science and ophthalmology. Their success in using elegant experiments to answer fundamental questions about vision demonstrates the power of basic research to improve human health. Their example has motivated generations of scientists to pursue careers in visual science and ophthalmology.

\subsection{Recognition and Honors}

In addition to the Nobel Prize, Granit, Hartline, and Wald received numerous other honors and awards for their work. Granit received the Order of the Polar Star in 1967 and was elected to the Royal Swedish Academy of Sciences. Hartline received the National Medal of Science in 1973 and was elected to the National Academy of Sciences. Wald received the National Medal of Science in 1967 and was elected to the National Academy of Sciences.

Their contributions to visual science and medicine are remembered and celebrated by researchers and clinicians around the world. Their discoveries continue to influence our understanding of vision and have led to numerous advances in the diagnosis and treatment of eye diseases, ensuring that their legacy will endure for generations to come.


\section{Modern Developments}

Granit, Hartline, and Wald's work on vision has led to modern ophthalmology. Understanding of phototransduction has enabled development of treatments for retinal diseases (anti-VEGF for macular degeneration, gene therapy for RPE65 mutations). Optogenetics uses light-sensitive proteins to restore vision in blind retinas. Understanding of retinal circuitry has informed development of retinal prostheses (Argus II) and optogenetic vision restoration therapies.


\section{Bibliography}

\begin{thebibliography}{9}
\bibitem{nobel-1967}
Nobel Prize Outreach, ``The Nobel Prize in Physiology or Medicine 1967,''\emph{NobelPrize.org}.\\
\url{https://www.nobelprize.org/prizes/medicine/1967/summary/}

\bibitem{lecture-1967}
Ragnar Granit, ``Nobel Lecture,''\emph{NobelPrize.org}. Winner-authored primary source; downloadable from the official Nobel Prize archive.\\
\url{https://www.nobelprize.org/prizes/medicine/1967/granit/lecture/}
\end{thebibliography}
